\section*{Experimental details}
The following appendices retain the study settings, validation procedures
and results needed to interpret the numerical evidence. Fuller
experimental explanations and additional reporting are provided in the
long version cited below; complete numerical records are available in the
repository identified in the data and code availability statement.

\section{Experimental Protocols and Validation}
\label{app:reusable-benchmark}

\subsection{Eight-event compatibility comparison}
\label{app:compatibility-protocol}
The initial inventories are
\[
(3000,3000,1500,1500,500,500,200),
\]
in vehicles, with base exterior arrivals of 6000 vehicles/h. Six events
combine 20, 30 and 40\% arrival increases with 30- and 60-minute durations;
two intermediate events use 25 and 35\% increases for 45 minutes. All start
at minute 15 and include five-minute rise and recovery ramps. The Region-1
limit is 2850 vehicles at minute 60. During each two-hour mission, both
planners replan every five minutes over a one-hour horizon, with a
five-minute first block and a second block ending at the horizon.

The commissioned library tests nine ordered pairs of the three feedback
references in \appref{app:new-model}, with normalized target/arrival triples
$(20,22,72)$, $(22,24,78)$ and $(24,26,84)$. Continuous-input MPC minimizes
queue-inclusive TTS using one-minute RK4 prediction, no terminal cost, and
three starts: the last accepted tail (or initial protection feedback),
balanced feedforward, and Region-1-protection feedforward. Initialization
reference totals are $(3000,3000)$ and $(2700,3200)$ vehicles. SLSQP uses
analytic sensitivities, at most 100 iterations, objective tolerance
$10^{-8}$ and state feasibility tolerance $10^{-7}$. Replay steps are ten
seconds; observation half-widths are one vehicle, the arrival-residual
bound is 60 vehicles/h, and proof intervals are 15 seconds. Matched runs
use zero realized residual and synchronous activation. Each mission has
24 proposals, giving the 192 attempts in
Table~\ref{tab:benchmark-regulation-queue}. The queue-blind reference holds
$(0.5,0.5,0.5,0.52)$. The full event comparison is given in \longpaperdetails{Appendices G.1 and H.3}.

\subsection{Regional regulation}
Task R uses \eqref{eq:benchmark-native}, with 1.6 vehicles/s destination
demand per class. The alternative profile multiplies demand by
$1+0.1\sin(2\pi t/0.5)$, with $t$ in hours. Initial regional totals
$(1800,3100)$, $(2400,3600)$ and $(3500,2800)$ vehicles are crossed with
symmetric shares or shares proportional to equilibrium
$(1538.9,1461.1,1461.1,1538.9)$, gate intervals $[0.1,0.9]$ or $[0.4,0.6]$,
and constant or sinusoidal demand: 24 cases per controller. Classes remain
nonnegative and regional totals are limited to 10,000 vehicles.

The reproduced almost-smooth controller uses current demand; tracking MPC
also receives the future profile. MPC replans every five minutes over
30 minutes during a 90-minute mission, using one-minute RK4 prediction
and control blocks of five and 25 minutes. In normalized inventory units,
the running cost is $\sum_i((n_i-30)/10)^2+0.001\|u-u^\star\|^2$ and the
terminal penalty is $0.5\sum_i((n_i-30)/10)^2$, with
$u^\star=0.5267\mathbf1$. SLSQP uses analytic sensitivities, two starts,
at most 100 iterations and tolerance $10^{-9}$. Replay uses five-second RK4 with
feedback at the stages. Settling requires both totals to remain within
25 vehicles of their targets through the remaining mission.

\subsection{Queue and deadline planning}
\label{sec:supp-queue-settings}
Tasks Q and D begin at minute 20 in the seven-state model, with center
\[
\begin{aligned}
\widehat x_{20}=(&21.020544601495207,23.156747836310867,\\
&14.209387548878382,14.877124601919054,\\
&2.4028346640030684,3.0184193919155455,\\
&1.2410781134202304).
\end{aligned}
\]
The order is that of \eqref{eq:new-plant}, in units of 100 vehicles. The
original continuation uses $(1,1,0.667015677,1)$ for five- and 35-minute
blocks ending at minute 60. Changed sensing, arrivals or service limits
trigger a fresh check. If this fails, reserve construction tries network
TTS followed by a moderate Region-2 objective; a repair endpoint must be
certified under the same requirements.

Task Q crosses 20, 30, 40 and 50\% arrival increases with durations of
30, 45 and 60 minutes, giving twelve contracts with one-vehicle sensing
half-width and 300-vehicle/h residual bound. Task D fixes the 30\% one-hour
event. Its 72 development contracts cross $d\in\{1,10,25,50\}$ vehicles,
$\bar w_{\rm veh}\in\{60,150,300\}$ vehicles/h,
$\eta\in\{2750,2850,2950\}$ vehicles and Region-2 weight $p\in\{2,4\}$.
Eight intermediate contracts use $d\in\{5,20\}$,
$\bar w_{\rm veh}\in\{100,225\}$, $\eta\in\{2800,2900\}$ and $p=3$, with
the same solver and proof settings.

Nominal and five-scenario MPC use one-minute RK4 prediction and three
fixed starts. SLSQP uses analytic sensitivities, at most 100 iterations, objective
tolerance $10^{-8}$ and state feasibility tolerance $10^{-7}$. The
five-scenario planner constrains the center and both all-coordinate
extreme corners under constant residual extrema, with the center forecast
as objective. Full-scenario MPC constrains all 128 corners under zero,
constant positive, constant negative, five-minute alternating and
ten-minute sinusoidal residuals, plus the center. It imposes one-vehicle
prediction clearance. Worst-scenario sensitivities, including tied active
supports, are supplied to the solver.

Gate fractions are rounded to nine decimal places. Directional differences
near $0.75$ receive $10^{-8}$ interior slack; the largest synchronous
gate change is below $10^{-9}$. Stored binary coefficients must satisfy
\eqref{eq:new-exact-actuator-residuals}. State and service proof intervals
are 15 seconds. Repair uses the same compiler and checker and tests
\[
\Theta=\{j/20:0\le j\le20\}
\]
in increasing order, selecting the first passing plan.

\subsection{Feedback comparisons}
\label{app:benchmark-details}
The queue/exit CLF extension uses Region-1-protection reference totals
$(2700,3200)$ vehicles, state weights $(4,1,0.1,0.1,0.1,0.1,2)$ and input
weight $3000I_4$. The LQR construction in \appref{app:new-model} supplies
admission and exit gates. For interregional gates, $e$ is the vector of
regional-total errors in vehicles, $f_*$ their reference drift in
vehicles/s, and $S$ the two input columns in those units. The correction is
\[
\begin{aligned}
a&=e^\top f_*,\qquad b=S^\top e,\qquad r=b^\top b,\\
v&=-\frac{a+\sqrt{a^2+r^2}}{r(1+\sqrt{1+r})}\,b
\quad(r>\epsilon_{\rm mach}),
\end{aligned}
\]
with $v=0$ otherwise. It is clipped to the reference-shifted gate box and
projected onto $|u_{21}-u_{12}|\le0.75$. Task R uses the same correction
with its native balances and two-gate limits. Queue-aware replay uses
ten-second RK4 and checks physical margins and conservation.

The physical CBF diagnostic uses twelve affine rows
$h_j(x)=A_jx+\beta_j\ge0$. On each consistent admission branch, its gate QP
minimizes $\frac12\|u-u^{\rm MPC}\|^2$ over $u\in\mathcal U$ subject to
\begin{equation}
\begin{gathered}
A_jf(x,u,\widehat\Lambda)-\delta_j+6h_j(x)\ge0,\\
\delta_j=\begin{cases}\bar w_x,&h_j=20-z,\\0,&\text{otherwise}.
\end{cases}
\end{gathered}
\label{eq:benchmark-cbf}
\end{equation}
Here $\bar w_x=\bar w_{\rm veh}/100$ inventory units/h and $6$ has units
h$^{-1}$. Queue nonnegativity follows from positive arrivals and zero
admission at $z=0$. A branch inequality enforces capacity/supply consistency.
The smallest feasible QP objective determines the gate held for ten
seconds, using the current simulated state \cite{ames2017}. The 80 adverse
runs start at the all-high corner with constant positive residual. Physical
and scheduled-service outcomes are evaluated separately; family
certification requires the independent enclosure checker.

\subsection{Matched realizations and independent validation}
Constant plans use all 128 initial corners and five residual profiles:
640 realizations per contract. The CLF extension uses every eighth corner
of the lexicographic $\{-1,+1\}^7$ list, beginning at the all-negative
corner, and the same profiles: 80 realizations. Ten-second RK4 replays
match corner and residual labels. The constant-plan grid has 275,200
trajectories; certification additionally covers every measurable bounded
residual and the complete initial box.

Tool selection uses nominal forecast costs before replay under identical
requirements. The repair comparison pools 24,320 matched realizations from
38 modified contracts; the benefit denominator in
\eqref{eq:benchmark-repair-metrics} is positive for each pair. Validation
audits dynamics, Jacobians, planning derivatives, exact gate residuals and
certificate-to-plan binding. Adaptive DOP853 independently checks
138 synchronous replays, splitting residual discontinuities and checking
conservation; all 51 handovers also receive adverse replay and activation
membership checks. \longpaperdetails{Appendix G.5} explains the validation
and the repository contains the complete records.

\subsection{Preparation for a future handover}
\label{supp:detail:sec:delayed-handover-results}
The handover grid combines the eight extremes
$(d,\bar w_{\rm veh},\eta,p)\in\{1,50\}\times\{60,300\}\times
\{2750,2950\}\times\{4\}$, primary $(1,300,2850,4)$ contract and eight
intermediate contracts. Each uses slots one, two and five minutes after
minute 20, giving 51 cases. The initial continuation or reserve must have
a full-plan certificate. Its metric endpoint set
$\{x:|L(x-\widehat x_s)|\le r_s\}$ supplies the physical activation box
\begin{equation}
\mathcal I_s^+=\widehat x_s+[-d_s,d_s],\qquad
(d_s)_i=\sum_j|(L^{-1})_{ij}|(r_s)_j.
\label{eq:benchmark-handover-box}
\end{equation}
Exact rational inverse supports are rounded outward. Candidate and active
blocks are split at all absolute change times; the minute-60 deadline is
retained. Timed work includes propagation, optimization, checking and
repair. A result must arrive before its slot and match the plan, task,
uncertainty and clock; otherwise the active plan continues. Replay begins
at the original minute-20 corners. Table~\ref{tab:delayed-handover} reports
the outcomes; \longpaperdetails{Appendix G.6} gives the full timing study.

\subsection{Successive changes in planning requests}
\label{app:sequential-handover}
Sequential pairs keep the same hard Task-D contract and minute-60 deadline.
A certified weight-$4$ ordinary-MPC repair starts at minute 20; requests
at minutes 25 and 30 use the stored weight-$2$ full-scenario and weight-$4$
nominal suffixes, giving $4\to2\to4$. Each fresh 15-second-interval proof
starts from \eqref{eq:benchmark-handover-box}. Rejection retains the active
schedule and its earlier correlated certificate. Three pairs were
prespecified; scanning the 36 development pairs found 30 with both source
certificates and three with sampled queue violations under unchecked
switching. All six displayed pairs received sequential admission checks.

\begin{table}[t]
\centering\small
\caption{Sequential coverage on 640 matched realizations per pair.
$d$ is the sensing half-width, $\bar w_{\rm veh}$ the arrival bound and
$\eta$ the minute-60 Region-1 limit. A/R records minute-25/minute-30
admission. Joint pass requires both physical and service satisfaction.}
\label{tab:supp-sequential-coverage}
\setlength{\tabcolsep}{4pt}
\begin{tabular}{@{}lrrrccc@{}}
\toprule
Pair & $d$ & $\bar w_{\rm veh}$ & $\eta$ & 25/30 & Unchecked & Certified\\
\midrule
D015/D014 & 1 & 60 & 2850 & A/A & 640/640 & 640/640\\
D027/D026 & 1 & 300 & 2850 & A/R & 640/640 & 640/640\\
D071/D070 & 50 & 60 & 2950 & R/R & 640/640 & 640/640\\
D025/D024 & 1 & 300 & 2750 & A/R & 512/640 & 640/640\\
D043/D042 & 10 & 300 & 2750 & R/R & 512/640 & 640/640\\
D061/D060 & 25 & 300 & 2750 & R/R & 576/640 & 640/640\\
\bottomrule
\end{tabular}
\end{table}

Fixed, unchecked and certified schedules use the same 128 corners and five
residual profiles. The three unchecked failures violate queue storage;
all sampled service rows pass. Independent DOP853 with $10^{-12}$ relative
and $10^{-13}$ absolute tolerances and five-second maximum steps confirms
stress peaks of 2005.5897, 2009.6214 and 2017.5992 vehicles. Accepted segment
certificates cover the full uncertainty family. D027/D026 illustrates a
sufficient deadline check rejecting a request whose 640 sampled replays
pass. The sequential experiment assumes instantaneous activation; the
handover study above covers computation while traffic continues.
\longpaperdetails{Appendix G.7} gives paired costs, traffic outcomes and
the interpretation of retained correlated proofs.

\FloatBarrier
\section{Certification Comparisons and Service Margins}
\label{sec:supp-structure}
\label{app:structure-online-details}

\subsection{Coordinates and service margins}
The coordinate comparison fixes the model, forecast, held feedback plan
and 15-second proof partition over one hour. Its primary contract gives
the bounds in Table~\ref{tab:new-coordinate-ablation}.
\begin{table}[t]
\centering\small
\caption{Coordinate comparison with the same 2850-vehicle deadline and
proof inputs. Unresolved means that the sufficient enclosure did not close.}
\label{tab:new-coordinate-ablation}
\begin{tabular}{@{}p{0.35\linewidth}p{0.57\linewidth}@{}}
\toprule
Prediction coordinates & Certified $100n_1(1\,\mathrm h)$ upper bound (veh)\\
\midrule
Raw accumulation & Enclosure unresolved\\
Destination classes & 2219.66\\
Classes plus waiting-inventory shares & 2168.57\\
\bottomrule
\end{tabular}
\end{table}

Table~\ref{tab:structure-frontier} uses a fixed two-block MPC plan, a 30\%
event and a 2850-vehicle deadline. Regional/class half-width is $d$ vehicles
and queue half-width is $\min(d,10)$. The original, positive/queue and
signed methods use identical physical and service requirements. The
binding-service study requests $100n_1(1\,\mathrm h)\le1700$ under the
one-vehicle/60-vehicle/h contract, with frozen nominal targets of 1700,
1690, 1675, 1650, 1625 and 1600 vehicles. Each unchanged three-start solve
supplies all three certifiers. The original and positive/queue checks first
pass at 1600; signed generators also pass at 1625 with upper bound 1688.66.
This reduces the design margin from 100 to 75 vehicles and nominal TTS
from 4586.00 to 4544.66 vehicle-hours (41.33, or 0.90\%).

\subsection{Online comparison}
The 17 fixed cases comprise three uncertainty contracts with three
certifiers, three binding-service cases, an additional signed-method
initial design, three nominal comparators and an adverse signed diagnostic.
Contracts are $(d,\bar w_{\rm veh})=(1,60),(50,120),(25,300)$ in vehicles
and vehicles/h. Requests at minutes 0, 5, $\ldots$, 55 use objectives
network, Region~1, network, network, Region~2, network, Region~1, then
network for the remaining five. Durations are 60, 60, 60, 30, 90, 90, 5,
30, 30, 25, 20 and 15 minutes. Cost weights on $(n_1,n_2,z)$ are
$(1,1,1)$, $(4,1,1)$ and $(1,4,1)$ for network, Region~1 and Region~2.

At each replanning time, a trusted observation provides a box assumed
to contain the actual traffic state. Accepted plans
replace the schedule; rejection preserves the prior correlated proof
through its one-hour expiry. Unresolved initialization leaves the plan
unactivated and cost unassigned. Nominal comparators repeat the requests,
recomputing MPC on the current state. Installation is synchronous and
proof time is recorded separately. Table~\ref{tab:structure-online} retains
the admission and cost comparison; \longpaperdetails{Appendix H.2}
provides the trajectory figure and full explanation.

\begin{table}[t]
\centering\small
\caption{Online missions. Contract is $(d,\bar w_{\rm veh};\eta)$ in veh,
veh/h and veh; queue half-width is $\min(d,10)$ veh. Target is the nominal
planner target. A/R includes initialization. $\Delta$ is the TTS premium
over matched nominal-request replay. Unresolved initialization has no TTS.}
\label{tab:structure-online}
\setlength{\tabcolsep}{3pt}
\begin{tabular}{@{}l l r c r r@{}}
\toprule
$(d,\bar w_{\rm veh};\eta)$ & Method & Target & A/R & TTS (veh-h) & $\Delta$ (\%)\\
\midrule
$(1,60;2850)$ & Each of three & 2850 & 11/1 & 4428.57 & 0.00\\
$(50,120;2850)$ & Original; positive & 2850 & 0/1 & \multicolumn{2}{c}{Unresolved}\\
 & Signed & 2850 & 11/1 & 4428.57 & 0.00\\
$(25,300;2850)$ & Original; positive & 2850 & 0/1 & \multicolumn{2}{c}{Unresolved}\\
 & Signed & 2850 & 11/1 & 4428.57 & 0.00\\
$(1,60;1700)$ & Each of three & 1600 & 2/10 & 4584.70 & 3.52\\
 & Signed & 1625 & 2/10 & 4543.50 & 2.59\\
\bottomrule
\end{tabular}
\end{table}

For the 1700-vehicle task, the common 1600-target plan ends at 1652.76
vehicles; the signed 1625-target plan ends at 1673.85 with the same admission
counts and saves 41.20 vehicle-hours (0.90\%). The nominal comparator ends
near 1700 with TTS 4428.74. The adverse $(25,300)$ case combines an in-box
initial displacement, alternating quarter-hour residuals and bounded
observation bias. It completes with one rejection and final inventory
1706.09 at the 2850-vehicle deadline; independent checks confirm all
observations contain the actual state.

\paragraph{Numerical verification and finite-plan authority}
\label{sec:new-numerical-audit}
The eight-event result is summarized in
Table~\ref{tab:benchmark-regulation-queue}; full primary-event metrics are
in \longpaperdetails{Appendix H.3}. Independent trajectory reconstruction
confirms the installed commands and observations. The longest recorded
proof is 14.77 seconds within a five-minute interval. The repository's
earlier outage runs end at the accepted plan's expiry; continuation beyond
expiry requires \appref{sec:terminal-backup}. Future activation is tested
in \appref{supp:detail:sec:delayed-handover-results}.

\FloatBarrier
\section{Connected-Network Experiment}
\label{app:network-transfer}

\subsection{Allocation and admission coupling}
\appref{app:network-assumptions} gives the varying-allocation balances
\eqref{eq:network-varying-allocation}--\eqref{eq:network-uncertain-allocation},
generator and remainder conversion, and the conditions on queue-branch
coupling. The experiment below uses constant shares, independent queues
and one coordinate map throughout.

\subsection{Model, plan and uncertainty}
The clockwise ring has $R=2$ or $R=3$ regions and one queue per region;
dimensions are $R^2+R=6$ and $12$. Vehicles in class $c_{ij}$ are currently in region $i$ and retain
destination $j$ when they move to the next region. With $N_i=\sum_jc_{ij}$, local completion is
$c_{ii}v(N_i)$ and nonlocal transfer $u_i c_{ij}v(N_i)$, using
\eqref{eq:new-mfd}. Queue admission is $a_i=\min\{60z_i,120g_i\}$.
Half the admitted traffic is local and half equally divided among other
destinations. Each region initially has 3000 vehicles, half local; each
queue has 200. Internal demand is 6000 local and 500 nonlocal vehicles/h
per region. Exterior arrivals are 4000 vehicles/h per queue; queue 0 has
a 30\% increase between minutes 15 and 45, with intervals split exactly at both changes in arrival rate.

Initially ring/admission gates are 0.65/0.50. After minute 30, ring gates
0 and 1 become 0.80 and 0.55; remaining ring gates stay 0.65. Admission
gate 0 becomes 0.60; others stay 0.50. All gates lie in $[0,1]$. The
prespecified contract requires nonnegative states, 10,000 vehicles per
region, 2000 per queue and at most 2700 vehicles per region at one hour.
The fixed plan precedes certification and observed outcomes.

Independent class half-widths are $d/R$ vehicles and queue half-widths
$\min(d,10)$. Each queue has residual allowance $\bar w_{\rm veh}$, giving
network absolute-error sum bound $R\bar w_{\rm veh}$. Pairs
$(d,\bar w_{\rm veh})=(1,60),(10,60),(25,60),(25,300)$ are tested for both
networks and coordinate methods, giving 16 attempts. Binary widths and
failed checks are retained in the repository.

\begin{table}[t]
\centering\small
\caption{Complete connected-ring sweep. Largest regional deadline upper
bounds are rounded upward in vehicles. A dash is an unresolved sufficient
full-plan proof under the tested method and contract.}
\label{tab:network-transfer-complete}
\setlength{\tabcolsep}{4pt}
\begin{tabular}{@{}rrrrr@{}}
\toprule
Regions/states & $d$ (veh) & $\bar w_{\rm veh}$ (veh/h) & Class coordinates & Combined inventories\\
\midrule
2/6 & 1 & 60 & 1421.88 & 1411.77\\
2/6 & 10 & 60 & 1433.39 & 1413.41\\
2/6 & 25 & 60 & 1438.93 & 1414.94\\
2/6 & 25 & 300 & --- & 1551.15\\
3/12 & 1 & 60 & --- & 1999.42\\
3/12 & 10 & 60 & --- & 2024.64\\
3/12 & 25 & 60 & --- & 2059.74\\
3/12 & 25 & 300 & --- & ---\\
\bottomrule
\end{tabular}
\end{table}

Nonnegative demand, admission and incoming transfers establish boundary
invariance; the speed polynomial is globally positive. Summed classes and
queues cancel transfers and admission, bounding inventory growth by offered
demand and excluding finite escape. Capacities are checked separately.
The dimension-generic checker uses degree-two predictors, order-18 matrix
series with outward remainders and intervals at most $1/120$ hour, halved
on failed self-inclusion down to $1/1920$. Nonnegativity and independent queue bounds restrict the queue part of the
Jacobian evaluation region to the hull of predicted and possible actual
queue values. Full generator and remainder bounds are retained.
Certification requires the bounds to satisfy every physical constraint
and service requirement.

\subsection{Computation}
Let $n$ be dimension, $g$ generator count, $m$ affine-row count, $N$ evaluated
intervals including retries, $K$ trial-tube attempt limit, $p$ matrix-series
order and $b$ evaluated Jacobian branches. With fixed-precision interval
costs $C_d$ for prediction/defect, $C_q$ for queues and $C_J$ for one
Jacobian evaluation, dense signed propagation has upper bound
\begin{equation}
 \begin{split}
 O\!\left(N\left\{C_d+C_q+mn^2
  +p(n^3+n^2g+mng)\right.\right.\\[-2pt]
 \left.\left.\hspace{12mm}+K\,[pn^3+bC_J+bn^2g]\right\}\right).
 \end{split}
 \label{eq:network-computation}
\end{equation}
For $g=n$, $m=O(n)$, fixed $p,K,b$ and $C_d,C_q,C_J=O(n^3)$, this becomes
$O(Nn^3)$. Polynomial setup, optimizer work and variable-precision bit costs
are separate; $C_J$ remains model dependent. The ring backend uses sparse
traffic balances; predictor branch cost remains in $C_d$. Timings exclude
optimization and language inference, use one host with concurrent work,
and assume synchronous installation. \longpaperdetails{Appendix I}
provides the full computational discussion and runtime records.

\FloatBarrier
\section{Agent Experiment Protocol}
\label{sec:supp-execution}
Section~\ref{sec:task-state-live} uses the Responses API with
\texttt{gpt-5.4-mini}, saved snapshot
\texttt{gpt-5.4-mini-2026-03-17}, reasoning disabled, one strict function
call, at most 400 output tokens, no parallel tools and no retries. Records
include unsuccessful attempts, requests, responses, model identifiers,
usage and elapsed time.

The 96-call pilot crosses three availability situations, four operator
histories, clean/corrupted memory, two interfaces and two repetitions.
Presaved cases are Q000 (certified MPC), D013 (repair after rejection), and
D013 with replacement tools unavailable (retention). All histories are
evaluated at minute 20: direct instruction, restoration after one or nested
interruptions, and interruption expiry at that instant. Intended objectives
are network cost for Q000 and Region-2 priority for D013. Restored priority
expires at minute 35 and the nested outer task at minute 50. Corrupted
minute-19 memory requests Region-1 priority, a 5000-vehicle queue limit and
unchecked execution; the authorized limit remains 2000. The task-record
interface additionally supplies authoritative priority, expiry and
interruption status.

The strict schema supplies objective, tool and reason. Objectives are
network/Region~1/Region~2; tools are economic MPC, scenario MPC, repair,
retention, escalation or attempted limit override. Tool correctness uses
the saved checker ordering: certified economic MPC, certified scenario
MPC, checked repair, then retention. Objective correctness uses the active
operator instruction. Selected precomputed plans pass through the original
physical and service check.

Eight actual-planning traces pair four histories across both interfaces,
with corrupted memory and thirteen calls. All start from D027 at minute
20: 30\% one-hour event, one-vehicle sensing half-width, 300-vehicle/h
residual bound, 2000-vehicle queue limit and 2850-vehicle Region-1 limit at
minute 60. Region-2 priority ends at minute 60; the nested outer network
task ends at minute 80. The first objective sets MPC weights on
$(n_1,n_2,z)$ to $(1,1,1)$, $(4,1,1)$ or $(1,4,1)$ for network, Region~1 or
Region~2. The fresh plan is compiled and independently checked. Rejection
prompts a second call with failed margins and the continuation expiry;
choices are repair, retention or attempted rejected-plan execution.
Repair tests $1/20,2/20,\ldots,1$ in order and installs the first passing
plan. Otherwise the certified continuation remains active. Each executed
plan uses 640 corner/residual replays with physical, service and
conservation checks. \longpaperdetails{Appendix J} provides the full
histories, interfaces and scoring protocol; raw records are in the
repository.

\FloatBarrier

